% ch4.tex -- Module 6 ESSENTIALS, Chapter 4: Transformers and Power Transmission
% Knowledge points only (no questions). Condensed from src/mod6/ch4.tex
% (final full book): same notation and facts. The two figures are TikZ, copied from the
% full chapter (shrunk). Dot points: M6.IQ3.03, M6.IQ3.04, M6.IQ3.05.
\chapter{Transformers and Power Transmission}\label{ch4:chap}
\covers{M6.IQ3.03, M6.IQ3.04, M6.IQ3.05}

% chapter-local colours (prefix chFour, as in the full book)
\colorlet{chFourcopper}{hscseriesB!80!black}  % coil wire
\colorlet{chFouriron}{hscmuted!28}            % iron core
\colorlet{chFourcurrent}{hscseriesB}          % induced (eddy) currents

\begin{summarybox}{In this chapter}
\begin{itemize}
  \item How a transformer changes an AC voltage, and the ideal-transformer equations.
  \item Why real transformers waste energy, and how engineers reduce the losses.
  \item Why the grid uses high voltage, and other uses of transformers.
\end{itemize}
\end{summarybox}

%% ==================================================================
\section{The ideal transformer}\label{ch4:sec:ideal}
\covers{M6.IQ3.03}

\begin{bigidea}
A transformer changes the size of an AC voltage. A changing current in one coil makes a
changing flux in an iron core, and this induces an emf in a second coil. Energy is
conserved: voltage up means current down.
\end{bigidea}

\begin{keywords}
  \kw{Alternating current (AC)}{current that keeps reversing (DC flows one way only).}
  \kw{Primary / secondary coil}{the input coil (joined to the supply) / the output coil
    (joined to the load, the device that uses the energy).}
  \kw{Mutual induction}{an emf made in one coil by a changing current in a nearby coil.}
  \kw{Step-up / step-down}{more / fewer secondary turns: bigger / smaller voltage.}
\end{keywords}

\subsection{How a transformer works}

Both coils are wrapped around one closed \term{core}: a ring of \term{soft iron} (iron that
is easy to magnetise and demagnetise), as in \cref{ch4:fig:transformer}.

\textbf{How the output voltage is made (cause $\to$ effect):}
\begin{enumerate}
  \item AC in the primary makes a magnetic field that keeps changing size and direction.
  \item The iron core carries this changing flux to the secondary coil.
  \item The flux through each secondary turn changes, so an emf is induced in every turn
        (Faraday's law, $\varepsilon = -N\dfrac{\Delta\Phi}{\Delta t}$).
  \item The turns are in series, so their emfs add up: more turns, bigger voltage.
  \item The flux reverses with the supply, so the output is AC of the \emph{same frequency}.
\end{enumerate}

\begin{sidebyside}[0.58]
  \centering
  \begin{tikzpicture}[x=0.7cm, y=0.75cm]
    % iron core (closed loop)
    \fill[chFouriron, even odd rule] (0,0) rectangle (5.6,3.6) (1.1,0.9) rectangle (4.5,2.7);
    \draw[hscink!70, line width=0.6pt] (0,0) rectangle (5.6,3.6) (1.1,0.9) rectangle (4.5,2.7);
    % flux path (clockwise at this instant)
    \draw[hscfield, line width=0.8pt, dashed] (0.55,0.45) rectangle (5.05,3.15);
    \draw[field] (1.4,3.15) -- (4.2,3.15);
    \draw[field] (4.2,0.45) -- (1.4,0.45);
    \node[diagramlabel, text=hscfield, align=center] at (2.8,1.8)
      {changing\\ flux $\Phi$};
    % primary coil: 4 turns on the left limb
    \foreach \y in {1.2,1.6,2.0,2.4}
      \draw[chFourcopper, line width=1.6pt] (-0.14,\y) -- (1.24,\y+0.18);
    % secondary coil: 8 turns on the right limb
    \foreach \y in {0.95,1.17,...,2.5}
      \draw[chFourcopper, line width=1.6pt] (5.74,\y) -- (4.36,\y+0.11);
    % primary circuit and AC supply
    \draw[hscink, line width=0.7pt] (-0.14,2.4) -- (-2.2,2.4) -- (-2.2,2.2);
    \draw[hscink, line width=0.7pt] (-0.14,1.2) -- (-2.2,1.2) -- (-2.2,1.4);
    \draw[hscink, line width=0.7pt, fill=white] (-2.2,1.8) circle (0.25cm);
    \node at (-2.2,1.8) {$\sim$};
    \node[diagramlabel, align=center, below] at (-2.2,1.1) {AC\\ supply};
    \node[diagramlabel] at (-1.2,1.8) {$V_\mathrm{p}$};
    \node[diagramlabel, above] at (-1.2,2.4) {$I_\mathrm{p}$};
    % secondary circuit and lamp
    \draw[hscink, line width=0.7pt] (5.74,2.49) -- (7.6,2.49) -- (7.6,2.15);
    \draw[hscink, line width=0.7pt] (5.74,0.95) -- (7.6,0.95) -- (7.6,1.45);
    \draw[hscink, line width=0.7pt, fill=white] (7.6,1.8) circle (0.21cm);
    \draw[hscink, line width=0.7pt] (7.6,1.8) ++(45:0.21cm) -- ++(225:0.42cm)
      (7.6,1.8) ++(135:0.21cm) -- ++(315:0.42cm);
    \node[diagramlabel, align=center, below] at (7.6,0.85) {lamp\\ (load)};
    \node[diagramlabel] at (6.7,1.75) {$V_\mathrm{s}$};
    \node[diagramlabel, above] at (6.7,2.49) {$I_\mathrm{s}$};
    % labels
    \node[diagramlabel] at (-0.4,4.1) {primary coil ($N_\mathrm{p}$ turns)};
    \draw[hscmuted, line width=0.5pt] (0.3,3.8) -- (0.3,2.52);
    \node[diagramlabel] at (6.0,4.1) {secondary coil ($N_\mathrm{s}$ turns)};
    \draw[hscmuted, line width=0.5pt] (5.3,3.8) -- (5.3,2.58);
    \node[diagramlabel, below] at (2.8,-0.05) {soft-iron core (a closed loop)};
  \end{tikzpicture}
  \captionof{figure}{A step-up transformer ($N_\mathrm{s} > N_\mathrm{p}$). The flux arrows
    show one instant; they reverse every half-cycle. No wire joins the coils: no charge
    crosses, and the changing field carries the energy.}
  \label{ch4:fig:transformer}
\nextside
  \begin{itemize}
    \item \textbf{Why AC is needed.} Steady DC makes a \emph{steady} field, so
          $\Delta\Phi/\Delta t = 0$: no emf. DC gives only a short pulse when it is
          switched on or off.
    \item \textbf{Why $V \propto N$.} The same flux passes through every turn of both
          coils, so each turn has the same emf. So the voltage ratio equals the
          \term{turns ratio} $N_\mathrm{p} : N_\mathrm{s}$.
  \end{itemize}
\end{sidebyside}

\begin{formulacard}{The transformer (turns) equation}
  \cardformula{\frac{V_\mathrm{p}}{V_\mathrm{s}} = \frac{N_\mathrm{p}}{N_\mathrm{s}}}
  \begin{symbols}
    \sym{$V_\mathrm{p}$, $V_\mathrm{s}$}{voltage of the primary (input), of the secondary}{\si{V}}
    \sym{$N_\mathrm{p}$, $N_\mathrm{s}$}{number of turns on the primary, on the secondary}{no unit}
  \end{symbols}
  \whentouse{a transformer on AC (exact if ideal). Not for DC, not for currents.}
  \trap{primary values on top of both fractions. The primary is joined to the supply.}
\end{formulacard}

\textbf{What the output looks like.}
\begin{itemize}
  \item The frequency stays the same. The \term{peak} (largest) voltage is multiplied by
        $N_\mathrm{s}/N_\mathrm{p}$.
  \item Mains: \qty{230}{V} AC, \qty{50}{Hz} ($T = 1/f = \qty{20}{ms}$). \qty{230}{V} is an
        \term{effective value} (the steady voltage that heats a heater by the same amount);
        the peak is about \qty{325}{V}. Never mix peak and effective values.
  \item \textbf{A graph of the primary \emph{current}:} the flux follows the current, so
        the secondary emf follows the \term{gradient} (steepness) of the current graph.
        Flat: zero emf. Steepest: biggest emf. A straight-line rise: a constant emf. A sine
        current: a cosine-shaped emf, shifted by a quarter of a period.
\end{itemize}

\subsection{Power and current}

In an \term{ideal transformer} no energy is wasted, so power in equals power out. The
\term{power} (energy each second) of each coil is $P = VI$ (on the sheet).

\begin{formulacard}{Power in an ideal transformer}
  \cardformula{V_\mathrm{p} I_\mathrm{p} = V_\mathrm{s} I_\mathrm{s}}
  \begin{symbols}
    \sym{$V_\mathrm{p}$, $I_\mathrm{p}$}{voltage and current of the primary (input)}{\si{V}, \si{A}}
    \sym{$V_\mathrm{s}$, $I_\mathrm{s}$}{voltage and current of the secondary (output)}{\si{V}, \si{A}}
  \end{symbols}
  \whentouse{an ideal transformer, to find a current or a power.}
  \trap{step-up: \emph{voltage} up, current \emph{down}. Never multiply $V_\mathrm{p}$ by
    $I_\mathrm{s}$.}
\end{formulacard}

\begin{itemize}
  \item With the turns equation:
        $\dfrac{I_\mathrm{s}}{I_\mathrm{p}} = \dfrac{V_\mathrm{p}}{V_\mathrm{s}}
        = \dfrac{N_\mathrm{p}}{N_\mathrm{s}}$ (not printed in this form). The current ratio
        is upside down: step-down gives a \emph{bigger} secondary current.
  \item \textbf{The load decides the current.} The load takes the current it needs; the
        primary then draws just enough to supply that power.
\end{itemize}

%% ==================================================================
\section{Real transformers: energy losses and efficiency}\label{ch4:sec:losses}
\covers{M6.IQ3.04}

\begin{bigidea}
Real transformers waste a little energy, mostly as heat. So the power out is a bit less
than the power in. A closed iron core, thick wire and a laminated core keep the losses
small.
\end{bigidea}

\begin{keywords}
  \kw{Incomplete flux linkage}{some flux from the primary misses the secondary.}
  \kw{Eddy currents}{loops of current induced inside a solid piece of metal by a changing
    magnetic flux.}
  \kw{Laminated core}{a core made of thin iron sheets with insulation between them.}
\end{keywords}

Real: $V_\mathrm{s}I_\mathrm{s} < V_\mathrm{p}I_\mathrm{p}$. The lost energy becomes heat
(and a little sound). \Cref{ch4:tab:losses} gives the three syllabus limitations.

\begin{table}
  \caption{The three main limitations of the ideal transformer model.}\label{ch4:tab:losses}
  \centering\small
  \begin{tabular}{@{}L{25mm}L{64mm}L{48mm}@{}}
    \toprule
    \colhead{Limitation} & \colhead{Cause and effect} & \colhead{How to reduce it}\\
    \midrule
    Incomplete flux linkage & Some primary flux leaks into the air and misses the
      secondary, so the secondary emf is smaller than predicted. &
      One \emph{closed} loop of soft iron for both coils. Wind the secondary directly over
      the primary.\\
    \addlinespace
    Resistive heating in the coils & The current heats the long copper coils ($I^2R$).
      This \term{resistive heating} is heat made when a current flows through a
      resistance. &
      Thick copper wire (thickest on the coil with the larger current). Cooling (vents,
      fans or oil) keeps the wire cool, and a cool wire has less resistance.\\
    \addlinespace
    Eddy currents in the core & The changing flux induces loops of current in the iron
      core, which heat it. & A laminated core, or a ferrite core.\\
    \bottomrule
  \end{tabular}
\end{table}

\textbf{Evaluating the ideal model.} Big transformers in power stations and
substations are about 99\% efficient; small ones often waste 10--20\%. So the ideal model
is a very good approximation, but never exact. Losses can be reduced, never removed.

\begin{formulacard}[notonsheet]{Efficiency}
  \cardformula{\eta = \frac{P_\mathrm{out}}{P_\mathrm{in}} \times 100\%
    = \frac{V_\mathrm{s}I_\mathrm{s}}{V_\mathrm{p}I_\mathrm{p}} \times 100\%}
  \begin{symbols}
    \sym{$\eta$}{efficiency (Greek ``eta'')}{no unit (\%)}
    \sym{$P_\mathrm{out}$, $P_\mathrm{in}$}{useful power out, power in}{\si{W}}
  \end{symbols}
  \whentouse{a real transformer (ideal: $\eta = 100\%$). Energy wasted:
    $\Delta E = (P_\mathrm{in} - P_\mathrm{out})\,\Delta t$ (from $P = \Delta E/\Delta t$).}
  \derive{a definition: useful power out $\div$ power in.}
  \trap{the efficiency goes on the \emph{power}, not on the voltage. Use $0.92$, not $92$,
    inside a formula.}
\end{formulacard}

\subsection{Eddy currents and the laminated core}

\begin{itemize}
  \item The core is a conductor with a changing flux. So eddy currents flow in it, in loops
        at right angles to the flux. By Lenz's law they oppose the change in flux. They heat the core.
  \item A laminated core is thin iron sheets, each with a thin insulating layer. The sheets
        lie \emph{along} the flux, so the flux passes easily. The insulation cuts the
        eddy-current paths: each loop stays inside one sheet (\cref{ch4:fig:eddy}).
\end{itemize}

\begin{sidebyside}[0.6]
  \centering
  \begin{tikzpicture}[x=0.75cm, y=0.75cm]
    \tikzset{chFourdot/.pic={\draw[hscfield, line width=0.6pt, fill=white] (0,0) circle (0.12);
      \fill[hscfield] (0,0) circle (0.035);}}
    % (a) solid core
    \fill[chFouriron] (0,0) rectangle (3,3);
    \draw[hscink!70, line width=0.6pt] (0,0) rectangle (3,3);
    \draw[chFourcurrent, line width=1.3pt, rounded corners=7pt] (0.5,0.5) rectangle (2.5,2.5);
    \draw[chFourcurrent, line width=1.3pt, -{Stealth[length=2.6mm]}] (1.3,2.5) -- (1.75,2.5);
    \draw[chFourcurrent, line width=1.3pt, -{Stealth[length=2.6mm]}] (1.7,0.5) -- (1.25,0.5);
    \draw[chFourcurrent, line width=1.3pt, -{Stealth[length=2.6mm]}] (2.5,1.7) -- (2.5,1.25);
    \draw[chFourcurrent, line width=1.3pt, -{Stealth[length=2.6mm]}] (0.5,1.3) -- (0.5,1.75);
    \foreach \p in {(1.5,1.5),(0.25,0.25),(2.75,0.25),(0.25,2.75),(2.75,2.75)}
      \pic at \p {chFourdot};
    \node[font=\small\bfseries, text=hscink] at (1.5,3.5) {(a) solid core};
    % (b) laminated core
    \begin{scope}[xshift=2.9cm]
      \fill[chFouriron] (0,0) rectangle (3,3);
      \foreach \x in {0.5,1.0,...,2.6}
        \draw[hscink, line width=1.3pt] (\x,0) -- (\x,3);
      \draw[hscink!70, line width=0.6pt] (0,0) rectangle (3,3);
      \foreach \x in {0,0.5,...,2.6} {
        \draw[chFourcurrent, line width=1pt, rounded corners=3pt]
          (\x+0.13,0.35) rectangle (\x+0.37,2.65);
        \draw[chFourcurrent, line width=1pt, -{Stealth[length=2mm]}]
          (\x+0.37,1.7) -- (\x+0.37,1.3);
        \draw[chFourcurrent, line width=1pt, -{Stealth[length=2mm]}]
          (\x+0.13,1.3) -- (\x+0.13,1.7);
      }
      \node[font=\small\bfseries, text=hscink] at (1.5,3.5) {(b) laminated core};
      \node[diagramlabel, anchor=west] (ins) at (3.3,2.4) {insulating layer};
      \draw[hscmuted, line width=0.4pt] (ins.west) -- (2.5,2.2);
      \node[diagramlabel, anchor=west] (sh) at (3.3,1.5) {thin iron sheet};
      \draw[hscmuted, line width=0.4pt] (sh.west) -- (2.8,1.5);
      \node[diagramlabel, anchor=west] (ed) at (3.3,0.6) {eddy current};
      \draw[hscmuted, line width=0.4pt] (ed.west) -- (2.87,0.6);
    \end{scope}
    \pic at (-0.45,-0.6) {chFourdot};
    \node[diagramlabel, anchor=west] at (-0.25,-0.6)
      {flux $\Phi$ points out of the page in both cores};
  \end{tikzpicture}
  \captionof{figure}{A slice through the core: (a) large eddy currents, (b) small ones. The
    flux comes out of the page and is growing. So the eddy currents flow clockwise (they
    reverse every half-cycle).}
  \label{ch4:fig:eddy}
\nextside
  \begin{itemize}
    \item \textbf{Why this works:} each thin loop surrounds only a thin strip of flux
          (small emf) and has a narrow, high-resistance path. So the eddy currents and
          their heat are much smaller.
    \item Small chargers use \term{ferrite} cores. Ferrite is a magnetic ceramic that
          conducts very poorly. Cooling does
          \emph{not} stop eddy currents.
  \end{itemize}
\end{sidebyside}


%% ==================================================================
\section{Sending electricity at high voltage}\label{ch4:sec:transmission}
\covers{M6.IQ3.05}

\begin{bigidea}
Power lines have resistance, so their current makes heat: $I^2R$. Stepping the voltage
up makes the current small, so much less energy is wasted.
\end{bigidea}


\begin{keywords}
  \kw{Transmission line}{a long, thick wire on tall towers that carries electricity far.}
  \kw{Voltage drop}{the voltage used up pushing current through a resistance, $IR$.}
  \kw{Substation}{a place with transformers that change the voltage.}
\end{keywords}

A long line has several ohms of resistance $R$. The current $I$ is the same all along it,
so the town gets only $V - IR$ (Ohm's law, \cref{ch4:fig:line}).

\begin{figure}
  \centering
  \begin{tikzpicture}[x=1cm, y=1cm]
    \draw[hscink, line width=0.7pt] (0,0.95) -- (0,1.6) -- (2.6,1.6);
    \draw[hscink, line width=0.7pt] (4.4,1.6) -- (7.4,1.6) -- (7.4,1.1);
    \draw[hscink, line width=0.7pt] (7.4,0.1) -- (7.4,-0.4) -- (0,-0.4) -- (0,0.25);
    \draw[hscink, line width=0.7pt, fill=white] (0,0.6) circle (0.35);
    \node at (0,0.6) {$\sim$};
    \draw[hscink, line width=0.7pt, fill=hscforce!10] (2.6,1.4) rectangle (4.4,1.8);
    \node[diagramlabel] at (3.5,1.6) {$R$};
    \draw[hscink, line width=0.7pt, fill=hscbrand!8] (7.0,0.1) rectangle (7.8,1.1);
    \node[diagramlabel, rotate=90] at (7.4,0.6) {town};
    \draw[velocity] (1.0,1.6) -- (1.9,1.6) node[midway, above, text=hscvel] {$I$};
    \node[diagramlabel, align=center, anchor=east] at (-0.45,0.6)
      {power station\\ sends voltage $V$};
    \node[diagramlabel, align=center, above] at (3.5,1.85)
      {line resistance $R$ (both wires)\\ voltage drop $V_\mathrm{drop} = IR$};
    \node[diagramlabel, align=center, anchor=west] at (7.95,0.6)
      {voltage left\\ for the town\\ $= V - IR$};
    \node[diagramlabel, text=hscforce] at (3.5,1.0) {heat: $P_\mathrm{loss} = I^2R$};
  \end{tikzpicture}
  \caption{A simple model of a transmission line. The current $I$ is the same all the way
    around the loop.}
  \label{ch4:fig:line}
\end{figure}

\textbf{Why high voltage (cause $\to$ effect):}
\begin{enumerate}
  \item A power station sends a fixed power. Since $P = VI$, a higher voltage needs a
        smaller current.
  \item The loss $I^2R$ depends on the \emph{square} of the current: 10 times the voltage
        gives $\tfrac{1}{10}$ of the current and $\tfrac{1}{100}$ of the loss. $R$ stays the
        same. (Thicker wires also help, but they are heavy and very expensive.)
\end{enumerate}

\begin{formulacard}[notonsheet, unbreakable]{Power lost in a transmission line}
  \cardformula{P_\mathrm{loss} = I^2R}
  \begin{symbols}
    \sym{$P_\mathrm{loss}$}{power turned into heat in the line}{\si{W}}
    \sym{$I$}{current in the line}{\si{A}}
    \sym{$R$}{total resistance of the line (both wires)}{\si{\ohm}}
  \end{symbols}
  \whentouse{heat lost in lines or coils. Find $I = P/V$ first; the town gets
    $P - P_\mathrm{loss}$.}
  \derive{$P = VI$ for the line only, with $V_\mathrm{drop} = IR$ (both on the sheet):
    $P_\mathrm{loss} = V_\mathrm{drop}\,I = (IR)\,I = I^2R$.}
  \trap{do not use $\dfrac{V^2}{R}$ with the line's voltage (for example \qty{330}{kV}).
    There, $V$ must be the voltage drop $IR$, not the voltage sent.}
\end{formulacard}

\textbf{The grid.} The \term{electricity grid} is the network of power stations, lines and transformers.
Typical NSW voltages: power station about \qty{20}{kV} $\to$ step-up $\to$ transmission
lines \qty{330}{kV} or \qty{500}{kV} $\to$ step-down at substations $\to$ \qty{132}{kV} or
\qty{66}{kV} $\to$ \qty{11}{kV} $\to$ step-down on a power pole $\to$ homes \qty{230}{V}.
\begin{itemize}
  \item \textbf{Step up} at the power station: a low current in the long lines, so little
        energy is lost.
  \item \textbf{Step down} near towns: high voltages are dangerous. They can spark (arc)
        across air and need tall towers and big insulators. Appliances are built for
        \qty{230}{V}.
  \item \textbf{Why AC.} Transformers need a changing current: AC, not steady DC. (A few
        undersea links, such as Basslink, use DC.)
  \item \textbf{Other uses.} Step-up: microwave oven (\qty{230}{V} $\to$ about \qty{2}{kV}).
        Step-down: welder (\qty{230}{V} $\to$ \qty{50}{V}: a very large current melts
        metal); phone charger (\qty{230}{V} $\to$ \qty{5}{V}, small ferrite core, then
        electronics make DC).
\end{itemize}


%% ==================================================================
\section{Key experiments}\label{ch4:sec:experiments}
\covers{M6.IQ3.03, M6.IQ3.04}

\begin{experiment}{Testing the transformer equations}
\aim{To test $V_\mathrm{p}/V_\mathrm{s} = N_\mathrm{p}/N_\mathrm{s}$ for a real
  transformer, and to find its efficiency.}
\method{A demountable (take-apart) kit: laminated U-shaped iron core, iron \emph{I-bar} to close the
  top, slide-on coils; low-voltage AC, AC meters. Primary 600 turns at \qty{6.0}{V}:
  measure $V_\mathrm{s}$ for 150 to 1200 secondary turns. With a load, read all four
  meters. Then lift off the I-bar.}
\results{$V_\mathrm{s}$ against $N_\mathrm{s}$: a straight line through the origin, about
  5\% below the ideal line. With the load, $\eta \approx 90\%$. Opening the core
  cut $V_\mathrm{s}$ (300 turns) from \qty{2.86}{V} to about \qty{0.9}{V}.}
\conclusion{$V_\mathrm{s} \propto N_\mathrm{s}$, as predicted. The small gap is mainly flux
  leakage; $P_\mathrm{out} < P_\mathrm{in}$ (heat). The ideal model is close, not exact.}
\significance{It supports the model behind every grid transformer, and shows why cores are
  closed loops.}
\end{experiment}

%% ==================================================================
\begin{summarybox}{Remember}
\begin{itemize}
  \item A transformer needs AC: changing primary current $\to$ changing flux in the iron
        core $\to$ induced emf in the secondary (Faraday's law). Steady DC gives no output.
  \item $\dfrac{V_\mathrm{p}}{V_\mathrm{s}} = \dfrac{N_\mathrm{p}}{N_\mathrm{s}}$: more
        secondary turns, bigger voltage. The frequency stays the same.
  \item Ideal: $V_\mathrm{p}I_\mathrm{p} = V_\mathrm{s}I_\mathrm{s}$, so voltage up means
        current down: $\dfrac{I_\mathrm{s}}{I_\mathrm{p}} = \dfrac{N_\mathrm{p}}{N_\mathrm{s}}$.
  \item Real losses: flux leakage, $I^2R$ heating in the coils, eddy currents in the core.
        Fixes: closed soft-iron core, thick copper wire, laminated core.
  \item Not on the sheet: $\eta = \dfrac{P_\mathrm{out}}{P_\mathrm{in}} \times 100\%$ and
        $P_\mathrm{loss} = I^2R$.
  \item Fixed power: higher voltage $\to$ smaller current $\to$ much less $I^2R$ loss.
        Step up at the power station; step down in stages to \qty{230}{V} for safety.
\end{itemize}
\end{summarybox}
